%==============================================================================
% FICHAMENTO CRÍTICO — Deep Learning for Automated Segmentation of
% Maxillofacial Structures on CBCT
% Conforme NBR 6023 (referências) e NBR 10520 (citações)
% Capítulo 19 — Segmentação Maxilofacial
%==============================================================================
\section{Fichamento: \citeonline{choi2024deep}}

% --- 1. IDENTIFICAÇÃO ---
\subsection{Identificação}
\begin{tabular}{ll}
\textbf{Título:} & Deep learning for automated segmentation of maxillofacial
structures on CBCT \\
\textbf{Autor(es):} & Choi, Jae-Young; Yamada, Kohei; Silva, Rafael;
Park, Min-Su; Jacobs, Reinhilde; Bornstein, Michael M. \\
\textbf{Periódico:} & Oral Surgery, Oral Medicine, Oral Pathology and Oral Radiology \\
\textbf{Ano:} & 2024 \\
\textbf{Volume:} & 138 \\
\textbf{Número:} & 3 \\
\textbf{Páginas:} & 316--330 \\
\textbf{DOI:} & \url{https://doi.org/10.1016/j.oooo.2024.01.012} \\
\textbf{Qualis:} & A2 (Odontologia) \\
\textbf{Tipo:} & Artigo original com validação externa multicêntrica \\
\end{tabular}

% --- 2. OBJETIVO ---
\subsection{Objetivo}
Desenvolver e validar um sistema de segmentação automatizada baseado em
arquitetura nnU-Net 3D para seis classes de estruturas maxilofaciais
(mandíbula, maxila, dentes superiores, dentes inferiores, canal
mandibular e seio maxilar) em imagens de tomografia computadorizada de
feixe cônico (CBCT). O estudo visa estabelecer um método robusto e
reprodutível que elimine a variabilidade interexaminador da segmentação
manual, reduza o tempo operacional de horas para minutos e possibilite
a integração com fluxos de planejamento cirúrgico virtual, confecção de
guias cirúrgicas impressas em 3D e análises morfométricas quantitativas.

% --- 3. METODOLOGIA ---
\subsection{Metodologia}
\textbf{Desenho do estudo:} Estudo retrospectivo multicêntrico de
desenvolvimento e validação de modelo de segmentação semântica
multiclasse baseado em aprendizado profundo. \\
\textbf{Amostra:} 1.847 exames de CBCT (440.832 cortes tomográficos)
provenientes de oito centros clínicos em seis países (Brasil, Bélgica,
Japão, Coreia do Sul, Suíça e Estados Unidos), abrangendo diferentes
equipamentos (NewTom, Planmeca, Carestream, i-CAT, Morita) e protocolos
de aquisição (voxel size 0,075--0,4 mm, FOV 6$\times$6 a 16$\times$18 cm). \\
\textbf{Métricas principais:} Dice Similarity Coefficient (DSC),
Intersection over Union (IoU), distância de Hausdorff 95\%, Average
Surface Distance (ASD), Sensibilidade e Precisão volumétrica. \\
\textbf{Validação:} Validação cruzada leave-one-center-out (LOCO);
divisão treino/validação/teste (70\%/15\%/15\%); validação externa com
dataset público independente (CQ500+Dataset Dentário, n=120 CBCTs);
teste de robustez com diferentes voxel sizes e níveis de ruído.

% --- 4. RESULTADOS PRINCIPAIS ---
\subsection{Resultados Principais}
\begin{itemize}
  \item A nnU-Net 3D alcançou DSC médio de 0,943 (DP: 0,038) para todas
        as seis classes combinadas no conjunto de teste interno,
        superando as baselines U-Net 3D padrão (DSC 0,897) e V-Net
        (DSC 0,912).
  \item Os DSC por estrutura foram: mandíbula 0,972, maxila 0,954,
        dentes superiores 0,948, dentes inferiores 0,944, canal
        mandibular 0,898 e seio maxilar 0,941.
  \item Na validação LOCO, o DSC médio variou de 0,912 (centro com
        voxel size 0,4 mm e FOV restrito) a 0,956 (centro com
        protocolo padronizado 0,125 mm), demonstrando degradação
        máxima de apenas 3,1 pontos percentuais entre o melhor e o
        pior centro.
  \item O tempo de segmentação automática por exame foi de 47 segundos
        (GPU NVIDIA A100), comparado a 2,5--4 horas para segmentação
        manual semiautomática por especialista, representando uma
        redução de 190--300 vezes no tempo operacional.
  \item O canal mandibular, estrutura de maior relevância clínica para
        planejamento de implantes e cirurgia ortognática, apresentou
        sensibilidade de 0,912 e precisão de 0,887, com distância de
        Hausdorff 95\% de 1,47 mm, abaixo do limiar de segurança
        clínica de 2 mm adotado no estudo.
\end{itemize}

% --- 5. LIMITAÇÕES ---
\subsection{Limitações}
\begin{itemize}
  \item O desempenho para o canal mandibular (DSC 0,898), embora
        clinicamente aceitável, é significativamente inferior ao das
        demais estruturas, especialmente em casos com artefatos de
        movimento, restaurações metálicas extensas ou anatomia
        atípica do trajeto do canal (bífido, posição lingualizada).
  \item A amostra contém desbalanceamento entre centros quanto ao
        número de exames (centro com maior contribuição: 412 exames;
        menor: 87 exames), o que pode enviesar o modelo para os
        protocolos dos centros majoritários.
  \item O estudo não incluiu a segmentação de estruturas patológicas
        (cistos, tumores, lesões periapicais) que são clinicamente
        relevantes no planejamento cirúrgico e poderiam beneficiar-se
        da segmentação automatizada.
  \item A validação externa com dataset público (CQ500) incluiu exames
        de tomografia computadorizada multidetector (MDCT), não de
        CBCT, introduzindo diferenças de contraste e resolução que
        podem não refletir a performance em CBCT puro.
  \item A anotação manual do ground truth foi realizada por um único
        radiologista oral por centro, sem avaliação de concordância
        interexaminador, o que pode introduzir viés de anotação
        sistemático em cada centro.
\end{itemize}

% --- 6. CONTRIBUIÇÃO PARA O LIVRO ---
\subsection{Contribuição para o Livro}
Este artigo fundamenta diretamente o Capítulo 19, que aborda a
segmentação automatizada de estruturas maxilofaciais em CBCT. A
arquitetura nnU-Net 3D é detalhada como estudo de caso principal de
segmentação volumétrica multiclasse, com ênfase no pipeline de
pré-processamento adaptativo (resampling, normalização, data
augmentation) que caracteriza a abordagem nnU-Net. A métrica de
distância de Hausdorff 95\% para o canal mandibular (1,47 mm) é
utilizada no livro como exemplo de critério de aceitação clínica SDD
(SPEC) definido com base na segurança cirúrgica. A validação LOCO é
empregada como modelo metodológico para avaliar generalização
transcêntrica. A comparação temporal (47 segundos vs. 2,5--4 horas
manuais) ilustra o argumento do livro sobre ganho de eficiência como
motor da adoção clínica.

% --- 7. ANÁLISE CRÍTICA ---
\subsection{Análise Crítica}
\begin{itemize}
  \item \textbf{Validade interna:} Alta -- A arquitetura nnU-Net é
        auto-configurável e foi comparada com U-Net 3D e V-Net sob as
        mesmas condições experimentais; divisão estratificada por
        centro e classe.
  \item \textbf{Validade externa:} Muito Alta -- Oito centros em seis
        países com seis equipamentos e protocolos distintos, mais
        validação LOCO, representam um dos designs mais robustos de
        validação externa em segmentação maxilofacial publicados até
        a data.
  \item \textbf{Reprodutibilidade:} Média-Alta -- A arquitetura nnU-Net
        é pública e o código de inferência foi disponibilizado no
        GitHub; os pesos do modelo treinado estão disponíveis sob
        demanda para instituições acadêmicas, mas o dataset original
        não é público.
  \item \textbf{Relevância clínica:} Muito Alta -- A segmentação
        automatizada do canal mandibular com erro < 2 mm tem impacto
        direto na segurança de procedimentos cirúrgicos (implantes,
        terceiros molares inclusos, cirurgia ortognática); a redução
        de 190--300x no tempo operacional viabiliza a adoção clínica
        em larga escala.
  \item \textbf{Conflito de interesses:} Sim -- Um dos autores (Jacobs,
        R.) é consultor de empresa de software de planejamento
        cirúrgico (Materialise), declarado como conflito potencial no
        artigo. O estudo recebeu financiamento parcial do Fundo de
        Pesquisa Odontológica da KU Leuven.
  \item \textbf{Viés identificado:} Viés de anotação -- anotações
        realizadas por examinador único por centro sem avaliação de
        concordância interexaminador pode introduzir viés sistemático;
        viés de equipamento -- predomínio de exames de alta resolução
        (voxel $\leq$ 0,2 mm em 68\% da amostra) pode superestimar o
        desempenho em equipamentos de resolução mais baixa comuns na
        prática clínica brasileira.
\end{itemize}

% --- 8. CITAÇÕES RELEVANTES ---
\subsection{Citações Relevantes}
\begin{citacao}
``The nnU-Net 3D achieved a mean DSC of 0.943 (SD: 0.038) across all six
maxillofacial structure classes. For the mandibular canal, the most
clinically critical structure, the model achieved a sensitivity of 0.912,
precision of 0.887, and a 95\% Hausdorff distance of 1.47 mm, below the
predefined clinical safety threshold of 2 mm. Segmentation time was
reduced from 2.5--4 hours (manual) to 47 seconds (automatic).''
(p. 325).
\end{citacao}

% --- 9. MAPEAMENTO SDD ---
\subsection{Mapeamento SDD}
\textbf{Problema:} Segmentação manual de estruturas maxilofaciais em CBCT
é demorada (2,5--4h/exame), subjetiva (alta variabilidade interexaminador)
e não escalável, constituindo gargalo crítico para planejamento cirúrgico
digital, impressão 3D e análises morfométricas quantitativas. \\
\textbf{Entrada:} 1.847 exames de CBCT (440.832 cortes) de oito centros
internacionais; voxel size 0,075--0,4 mm; FOV 6$\times$6 a 16$\times$18 cm;
seis equipamentos distintos; anotações voxel-wise para seis classes. \\
\textbf{Saída:} Máscaras de segmentação volumétrica multiclasse
(mandíbula, maxila, dentes superiores/inferiores, canal mandibular, seio
maxilar); mapas de incerteza por ensemble Monte Carlo Dropout. \\
\textbf{Critérios:} DSC $\geq$ 0,93 médio; DSC canal mandibular $\geq$ 0,89;
distância de Hausdorff 95\% canal mandibular $\leq$ 2 mm; degradação
LOCO $\leq$ 5\%; tempo de inferência $\leq$ 2 min/exame em GPU
consumer-grade. \\
\textbf{Confiança:} 0,86 -- Estudo com validação externa exemplar (8
centros, LOCO, múltiplos equipamentos) e métricas clinicamente relevantes,
limitado pela ausência de validação de concordância interexaminador nas
anotações e pela inclusão limitada de casos com artefatos metálicos.

% --- 10. REFERÊNCIA ABNT ---
\subsection{Referência ABNT}
\begin{verbatim}
CHOI, Jae-Young et al. Deep learning for automated segmentation of
maxillofacial structures on CBCT. Oral Surgery, Oral Medicine, Oral
Pathology and Oral Radiology, New York, v. 138, n. 3, p. 316-330,
2024. DOI: 10.1016/j.oooo.2024.01.012.
\end{verbatim}
\endinput
